\documentclass[12pt]{article}

\usepackage[margin=1in]{geometry}
\usepackage{amsmath}
\usepackage{amssymb}
\usepackage{bbm}
\usepackage{setspace}
\usepackage{color}
\definecolor{oucrimson}{RGB}{132,22,23}
\usepackage[colorlinks=true,
            urlcolor=oucrimson,
            linkcolor=black,
            citecolor=black]{hyperref}
%\doublespacing

% --------------------------------------------------
\title{Problem Set 4}
\author{Tristan Winkle}
\date{}

% --------------------------------------------------
\begin{document}
\maketitle
% --------------------------------------------------

% --------------------------------------------------
\section{Questions and Responses}
% --------------------------------------------------

\begin{enumerate} 
	\item What data sources might you scrape?
        \begin{itemize}
            \item[-] Developing web scraping skills has unlocked many new data opportunities for my research interests. A risk-free application involves scraping \href{https://reliefweb.int/}{ReliefWeb} to collect granular disaster declarations in developing countries. As a second option, I intend to attempt to scrape \href{https://www.clearancejobs.com/}{ClearanceJobs} to construct a dataset that captures human capital dynamics within mercenary labor markets.
        \end{itemize}

    \item[5(d).] What type of an object is mydf\$date?
        \begin{itemize}
            \item[-] Character
        \end{itemize}

    \item[6.7.] Verify that the two dataframe are different types: type class(df1) and class(df). What is the class of each? 
        \begin{itemize}
            \item[-] Gemini Pro 3.1 (since OSCER would not run sparklyr package):  The class of df1 is a local R dataframe (``tbl\_df"). The class of df is a "tbl\_spark".
        \end{itemize}

    \item[6.8.] Are the column names any different across the two objects? If so, why might that be?
        \begin{itemize}
            \item[-] They should be different. The iris df separates column names with ``." and Spark translates them to ``\_". 
        \end{itemize}
\end{enumerate}

% --------------------------------------------------
\end{document}
% --------------------------------------------------
